% ============================================================
% Gemeinsamer Preamble fuer alle 5 Open-Book-Spickzettel.
% Wird aus jedem klausur.tex via % ============================================================
% Gemeinsamer Preamble fuer alle 5 Open-Book-Spickzettel.
% Wird aus jedem klausur.tex via % ============================================================
% Gemeinsamer Preamble fuer alle 5 Open-Book-Spickzettel.
% Wird aus jedem klausur.tex via % ============================================================
% Gemeinsamer Preamble fuer alle 5 Open-Book-Spickzettel.
% Wird aus jedem klausur.tex via \input{../preamble.tex} geladen.
% Layout: A4 hochkant, 3 Spalten, 8pt, sehr kompakt.
% ============================================================

\usepackage[a4paper,margin=8mm,top=7mm,bottom=7mm,includefoot]{geometry}
\usepackage[utf8]{inputenc}
\usepackage[T1]{fontenc}
\usepackage[ngerman]{babel}
\usepackage{lmodern}
\usepackage{microtype}

\usepackage{amsmath,amssymb,amsfonts,mathtools}
\usepackage{multicol}
\usepackage{enumitem}
\usepackage{array,tabularx,booktabs,makecell}
\usepackage[dvipsnames,table]{xcolor}
\usepackage[explicit]{titlesec}
\usepackage{etoolbox}
\usepackage{fancyhdr}
\usepackage{lastpage}

% ---------- Farben ----------
\definecolor{secblue}{HTML}{1f4e79}
\definecolor{subblue}{HTML}{2e75b6}
\definecolor{boxbg}{HTML}{f2f2f2}
\definecolor{accent}{HTML}{c00000}
\definecolor{rulegray}{HTML}{888888}

% ---------- Spalten ----------
\setlength{\columnsep}{3mm}
\setlength{\columnseprule}{0.3pt}
\def\columnseprulecolor{\color{rulegray}}

% ---------- Listen kompakt ----------
\setlist{nosep,leftmargin=1.1em,labelsep=0.35em,topsep=0.2ex,partopsep=0pt}
\setlist[itemize]{label=\textbullet}

% ---------- Abstaende ----------
\setlength{\parindent}{0pt}
\setlength{\parskip}{0.15ex}
\setlength{\abovedisplayskip}{0.3ex}
\setlength{\belowdisplayskip}{0.3ex}
\setlength{\abovedisplayshortskip}{0.2ex}
\setlength{\belowdisplayshortskip}{0.2ex}
\setlength{\jot}{1pt}

% ---------- Tabellen kompakt ----------
\renewcommand{\arraystretch}{1.05}
\setlength{\tabcolsep}{3pt}
\setlength{\arrayrulewidth}{0.3pt}
\arrayrulecolor{rulegray}

% ---------- Section-Stil (kein Vollblock; dick + farbige Linie darunter) ----------
\titleformat{\section}
  {\large\bfseries\color{secblue}}
  {}{0pt}{#1}
  [\vspace{-0.4ex}{\color{secblue}\hrule height 0.8pt}\vspace{0.1ex}]
\titlespacing*{\section}{0pt}{1.4ex plus 0.4ex}{0.5ex}

\titleformat{\subsection}
  {\small\bfseries\color{subblue}}
  {}{0pt}{#1}
\titlespacing*{\subsection}{0pt}{0.9ex plus 0.2ex}{0.2ex}

\titleformat{\subsubsection}
  {\small\bfseries}
  {}{0pt}{#1}
\titlespacing*{\subsubsection}{0pt}{0.5ex}{0.1ex}

% ---------- Mini-Helfer ----------

% \fact{...} Hervorgehobene Regel (gerahmt, hellgrau)
\newcommand{\fact}[1]{%
  \begingroup\setlength{\fboxsep}{2pt}%
  \colorbox{boxbg}{\parbox{\dimexpr\linewidth-2\fboxsep}{\small #1\strut}}%
  \endgroup\par\vspace{0.2ex}}

% \warn{...} Stolperfalle (roter Strich links)
\newcommand{\warn}[1]{%
  \begingroup\setlength{\fboxsep}{2pt}%
  \noindent{\color{accent}\rule{2pt}{\baselineskip}}\hspace{2pt}%
  \begin{minipage}[t]{\dimexpr\linewidth-6pt}\small #1\end{minipage}%
  \endgroup\par\vspace{0.2ex}}

% Math-Shortcuts
\newcommand{\R}{\mathbb{R}}
\newcommand{\N}{\mathbb{N}}
\newcommand{\Z}{\mathbb{Z}}
\newcommand{\Q}{\mathbb{Q}}
\newcommand{\C}{\mathbb{C}}
\DeclareMathOperator{\rg}{rg}
\DeclareMathOperator{\spur}{tr}
\DeclareMathOperator{\im}{im}
\DeclareMathOperator{\diag}{diag}
\DeclareMathOperator{\sign}{sgn}
\DeclareMathOperator*{\argmin}{arg\,min}
\DeclareMathOperator*{\argmax}{arg\,max}
\newcommand{\trans}{^{\!\top}}
\newcommand{\norm}[1]{\left\lVert #1 \right\rVert}
\newcommand{\abs}[1]{\left\lvert #1 \right\rvert}
\newcommand{\inner}[2]{\langle #1,\,#2\rangle}
\newcommand{\dx}{\,\mathrm{d}x}
\newcommand{\dy}{\,\mathrm{d}y}
\newcommand{\dt}{\,\mathrm{d}t}

% Zeigt eine Display-Formel zentriert; skaliert nur runter wenn zu breit.
\newlength{\fitw}
\newcommand{\fitline}[1]{%
  \par\nointerlineskip
  \settowidth{\fitw}{$\displaystyle #1$}%
  \ifdim\fitw>\linewidth
    \noindent\resizebox{\linewidth}{!}{$\displaystyle #1$}\par
  \else
    \[\displaystyle #1\]%
  \fi}

% Tabular fuer fixe Spaltenbreiten (Spickzettel-typisch)
\newenvironment{tightt}[1]
  {\par\noindent\begin{tabular}{#1}}
  {\end{tabular}\par}
% Fuer X-Spalten direkt \begin{tabularx}{\linewidth}{...} verwenden
% (kann nicht in \newenvironment gewrappt werden -- TX@get@body Konflikt)

% Globale Schutzschalter gegen Overflow
\emergencystretch=2em
\tolerance=2000
\hbadness=10000
\hfuzz=2pt
\vfuzz=2pt

% ---------- Kopf-/Fusszeile ----------
\pagestyle{fancy}
\fancyhf{}
\renewcommand{\headrulewidth}{0pt}
\renewcommand{\footrulewidth}{0pt}
\fancyfoot[L]{\scriptsize\color{rulegray}\textit{\sheettitle}}
\fancyfoot[R]{\scriptsize\color{rulegray}Seite \thepage\,/\,\pageref{LastPage}}
\setlength{\footskip}{8pt}

% Titel-Header oben auf Seite 1
\newcommand{\sheetheader}[2]{%
  {\Large\bfseries\color{secblue} #1}\hfill
  {\small\color{rulegray} #2}\par
  \vspace{0.4ex}{\color{secblue}\hrule height 0.7pt}\vspace{1ex}}

% Default-Titel (wird ueberschrieben)
\newcommand{\sheettitle}{Open-Book Spickzettel}
 geladen.
% Layout: A4 hochkant, 3 Spalten, 8pt, sehr kompakt.
% ============================================================

\usepackage[a4paper,margin=8mm,top=7mm,bottom=7mm,includefoot]{geometry}
\usepackage[utf8]{inputenc}
\usepackage[T1]{fontenc}
\usepackage[ngerman]{babel}
\usepackage{lmodern}
\usepackage{microtype}

\usepackage{amsmath,amssymb,amsfonts,mathtools}
\usepackage{multicol}
\usepackage{enumitem}
\usepackage{array,tabularx,booktabs,makecell}
\usepackage[dvipsnames,table]{xcolor}
\usepackage[explicit]{titlesec}
\usepackage{etoolbox}
\usepackage{fancyhdr}
\usepackage{lastpage}

% ---------- Farben ----------
\definecolor{secblue}{HTML}{1f4e79}
\definecolor{subblue}{HTML}{2e75b6}
\definecolor{boxbg}{HTML}{f2f2f2}
\definecolor{accent}{HTML}{c00000}
\definecolor{rulegray}{HTML}{888888}

% ---------- Spalten ----------
\setlength{\columnsep}{3mm}
\setlength{\columnseprule}{0.3pt}
\def\columnseprulecolor{\color{rulegray}}

% ---------- Listen kompakt ----------
\setlist{nosep,leftmargin=1.1em,labelsep=0.35em,topsep=0.2ex,partopsep=0pt}
\setlist[itemize]{label=\textbullet}

% ---------- Abstaende ----------
\setlength{\parindent}{0pt}
\setlength{\parskip}{0.15ex}
\setlength{\abovedisplayskip}{0.3ex}
\setlength{\belowdisplayskip}{0.3ex}
\setlength{\abovedisplayshortskip}{0.2ex}
\setlength{\belowdisplayshortskip}{0.2ex}
\setlength{\jot}{1pt}

% ---------- Tabellen kompakt ----------
\renewcommand{\arraystretch}{1.05}
\setlength{\tabcolsep}{3pt}
\setlength{\arrayrulewidth}{0.3pt}
\arrayrulecolor{rulegray}

% ---------- Section-Stil (kein Vollblock; dick + farbige Linie darunter) ----------
\titleformat{\section}
  {\large\bfseries\color{secblue}}
  {}{0pt}{#1}
  [\vspace{-0.4ex}{\color{secblue}\hrule height 0.8pt}\vspace{0.1ex}]
\titlespacing*{\section}{0pt}{1.4ex plus 0.4ex}{0.5ex}

\titleformat{\subsection}
  {\small\bfseries\color{subblue}}
  {}{0pt}{#1}
\titlespacing*{\subsection}{0pt}{0.9ex plus 0.2ex}{0.2ex}

\titleformat{\subsubsection}
  {\small\bfseries}
  {}{0pt}{#1}
\titlespacing*{\subsubsection}{0pt}{0.5ex}{0.1ex}

% ---------- Mini-Helfer ----------

% \fact{...} Hervorgehobene Regel (gerahmt, hellgrau)
\newcommand{\fact}[1]{%
  \begingroup\setlength{\fboxsep}{2pt}%
  \colorbox{boxbg}{\parbox{\dimexpr\linewidth-2\fboxsep}{\small #1\strut}}%
  \endgroup\par\vspace{0.2ex}}

% \warn{...} Stolperfalle (roter Strich links)
\newcommand{\warn}[1]{%
  \begingroup\setlength{\fboxsep}{2pt}%
  \noindent{\color{accent}\rule{2pt}{\baselineskip}}\hspace{2pt}%
  \begin{minipage}[t]{\dimexpr\linewidth-6pt}\small #1\end{minipage}%
  \endgroup\par\vspace{0.2ex}}

% Math-Shortcuts
\newcommand{\R}{\mathbb{R}}
\newcommand{\N}{\mathbb{N}}
\newcommand{\Z}{\mathbb{Z}}
\newcommand{\Q}{\mathbb{Q}}
\newcommand{\C}{\mathbb{C}}
\DeclareMathOperator{\rg}{rg}
\DeclareMathOperator{\spur}{tr}
\DeclareMathOperator{\im}{im}
\DeclareMathOperator{\diag}{diag}
\DeclareMathOperator{\sign}{sgn}
\DeclareMathOperator*{\argmin}{arg\,min}
\DeclareMathOperator*{\argmax}{arg\,max}
\newcommand{\trans}{^{\!\top}}
\newcommand{\norm}[1]{\left\lVert #1 \right\rVert}
\newcommand{\abs}[1]{\left\lvert #1 \right\rvert}
\newcommand{\inner}[2]{\langle #1,\,#2\rangle}
\newcommand{\dx}{\,\mathrm{d}x}
\newcommand{\dy}{\,\mathrm{d}y}
\newcommand{\dt}{\,\mathrm{d}t}

% Zeigt eine Display-Formel zentriert; skaliert nur runter wenn zu breit.
\newlength{\fitw}
\newcommand{\fitline}[1]{%
  \par\nointerlineskip
  \settowidth{\fitw}{$\displaystyle #1$}%
  \ifdim\fitw>\linewidth
    \noindent\resizebox{\linewidth}{!}{$\displaystyle #1$}\par
  \else
    \[\displaystyle #1\]%
  \fi}

% Tabular fuer fixe Spaltenbreiten (Spickzettel-typisch)
\newenvironment{tightt}[1]
  {\par\noindent\begin{tabular}{#1}}
  {\end{tabular}\par}
% Fuer X-Spalten direkt \begin{tabularx}{\linewidth}{...} verwenden
% (kann nicht in \newenvironment gewrappt werden -- TX@get@body Konflikt)

% Globale Schutzschalter gegen Overflow
\emergencystretch=2em
\tolerance=2000
\hbadness=10000
\hfuzz=2pt
\vfuzz=2pt

% ---------- Kopf-/Fusszeile ----------
\pagestyle{fancy}
\fancyhf{}
\renewcommand{\headrulewidth}{0pt}
\renewcommand{\footrulewidth}{0pt}
\fancyfoot[L]{\scriptsize\color{rulegray}\textit{\sheettitle}}
\fancyfoot[R]{\scriptsize\color{rulegray}Seite \thepage\,/\,\pageref{LastPage}}
\setlength{\footskip}{8pt}

% Titel-Header oben auf Seite 1
\newcommand{\sheetheader}[2]{%
  {\Large\bfseries\color{secblue} #1}\hfill
  {\small\color{rulegray} #2}\par
  \vspace{0.4ex}{\color{secblue}\hrule height 0.7pt}\vspace{1ex}}

% Default-Titel (wird ueberschrieben)
\newcommand{\sheettitle}{Open-Book Spickzettel}
 geladen.
% Layout: A4 hochkant, 3 Spalten, 8pt, sehr kompakt.
% ============================================================

\usepackage[a4paper,margin=8mm,top=7mm,bottom=7mm,includefoot]{geometry}
\usepackage[utf8]{inputenc}
\usepackage[T1]{fontenc}
\usepackage[ngerman]{babel}
\usepackage{lmodern}
\usepackage{microtype}

\usepackage{amsmath,amssymb,amsfonts,mathtools}
\usepackage{multicol}
\usepackage{enumitem}
\usepackage{array,tabularx,booktabs,makecell}
\usepackage[dvipsnames,table]{xcolor}
\usepackage[explicit]{titlesec}
\usepackage{etoolbox}
\usepackage{fancyhdr}
\usepackage{lastpage}

% ---------- Farben ----------
\definecolor{secblue}{HTML}{1f4e79}
\definecolor{subblue}{HTML}{2e75b6}
\definecolor{boxbg}{HTML}{f2f2f2}
\definecolor{accent}{HTML}{c00000}
\definecolor{rulegray}{HTML}{888888}

% ---------- Spalten ----------
\setlength{\columnsep}{3mm}
\setlength{\columnseprule}{0.3pt}
\def\columnseprulecolor{\color{rulegray}}

% ---------- Listen kompakt ----------
\setlist{nosep,leftmargin=1.1em,labelsep=0.35em,topsep=0.2ex,partopsep=0pt}
\setlist[itemize]{label=\textbullet}

% ---------- Abstaende ----------
\setlength{\parindent}{0pt}
\setlength{\parskip}{0.15ex}
\setlength{\abovedisplayskip}{0.3ex}
\setlength{\belowdisplayskip}{0.3ex}
\setlength{\abovedisplayshortskip}{0.2ex}
\setlength{\belowdisplayshortskip}{0.2ex}
\setlength{\jot}{1pt}

% ---------- Tabellen kompakt ----------
\renewcommand{\arraystretch}{1.05}
\setlength{\tabcolsep}{3pt}
\setlength{\arrayrulewidth}{0.3pt}
\arrayrulecolor{rulegray}

% ---------- Section-Stil (kein Vollblock; dick + farbige Linie darunter) ----------
\titleformat{\section}
  {\large\bfseries\color{secblue}}
  {}{0pt}{#1}
  [\vspace{-0.4ex}{\color{secblue}\hrule height 0.8pt}\vspace{0.1ex}]
\titlespacing*{\section}{0pt}{1.4ex plus 0.4ex}{0.5ex}

\titleformat{\subsection}
  {\small\bfseries\color{subblue}}
  {}{0pt}{#1}
\titlespacing*{\subsection}{0pt}{0.9ex plus 0.2ex}{0.2ex}

\titleformat{\subsubsection}
  {\small\bfseries}
  {}{0pt}{#1}
\titlespacing*{\subsubsection}{0pt}{0.5ex}{0.1ex}

% ---------- Mini-Helfer ----------

% \fact{...} Hervorgehobene Regel (gerahmt, hellgrau)
\newcommand{\fact}[1]{%
  \begingroup\setlength{\fboxsep}{2pt}%
  \colorbox{boxbg}{\parbox{\dimexpr\linewidth-2\fboxsep}{\small #1\strut}}%
  \endgroup\par\vspace{0.2ex}}

% \warn{...} Stolperfalle (roter Strich links)
\newcommand{\warn}[1]{%
  \begingroup\setlength{\fboxsep}{2pt}%
  \noindent{\color{accent}\rule{2pt}{\baselineskip}}\hspace{2pt}%
  \begin{minipage}[t]{\dimexpr\linewidth-6pt}\small #1\end{minipage}%
  \endgroup\par\vspace{0.2ex}}

% Math-Shortcuts
\newcommand{\R}{\mathbb{R}}
\newcommand{\N}{\mathbb{N}}
\newcommand{\Z}{\mathbb{Z}}
\newcommand{\Q}{\mathbb{Q}}
\newcommand{\C}{\mathbb{C}}
\DeclareMathOperator{\rg}{rg}
\DeclareMathOperator{\spur}{tr}
\DeclareMathOperator{\im}{im}
\DeclareMathOperator{\diag}{diag}
\DeclareMathOperator{\sign}{sgn}
\DeclareMathOperator*{\argmin}{arg\,min}
\DeclareMathOperator*{\argmax}{arg\,max}
\newcommand{\trans}{^{\!\top}}
\newcommand{\norm}[1]{\left\lVert #1 \right\rVert}
\newcommand{\abs}[1]{\left\lvert #1 \right\rvert}
\newcommand{\inner}[2]{\langle #1,\,#2\rangle}
\newcommand{\dx}{\,\mathrm{d}x}
\newcommand{\dy}{\,\mathrm{d}y}
\newcommand{\dt}{\,\mathrm{d}t}

% Zeigt eine Display-Formel zentriert; skaliert nur runter wenn zu breit.
\newlength{\fitw}
\newcommand{\fitline}[1]{%
  \par\nointerlineskip
  \settowidth{\fitw}{$\displaystyle #1$}%
  \ifdim\fitw>\linewidth
    \noindent\resizebox{\linewidth}{!}{$\displaystyle #1$}\par
  \else
    \[\displaystyle #1\]%
  \fi}

% Tabular fuer fixe Spaltenbreiten (Spickzettel-typisch)
\newenvironment{tightt}[1]
  {\par\noindent\begin{tabular}{#1}}
  {\end{tabular}\par}
% Fuer X-Spalten direkt \begin{tabularx}{\linewidth}{...} verwenden
% (kann nicht in \newenvironment gewrappt werden -- TX@get@body Konflikt)

% Globale Schutzschalter gegen Overflow
\emergencystretch=2em
\tolerance=2000
\hbadness=10000
\hfuzz=2pt
\vfuzz=2pt

% ---------- Kopf-/Fusszeile ----------
\pagestyle{fancy}
\fancyhf{}
\renewcommand{\headrulewidth}{0pt}
\renewcommand{\footrulewidth}{0pt}
\fancyfoot[L]{\scriptsize\color{rulegray}\textit{\sheettitle}}
\fancyfoot[R]{\scriptsize\color{rulegray}Seite \thepage\,/\,\pageref{LastPage}}
\setlength{\footskip}{8pt}

% Titel-Header oben auf Seite 1
\newcommand{\sheetheader}[2]{%
  {\Large\bfseries\color{secblue} #1}\hfill
  {\small\color{rulegray} #2}\par
  \vspace{0.4ex}{\color{secblue}\hrule height 0.7pt}\vspace{1ex}}

% Default-Titel (wird ueberschrieben)
\newcommand{\sheettitle}{Open-Book Spickzettel}
 geladen.
% Layout: A4 hochkant, 3 Spalten, 8pt, sehr kompakt.
% ============================================================

\usepackage[a4paper,margin=8mm,top=7mm,bottom=7mm,includefoot]{geometry}
\usepackage[utf8]{inputenc}
\usepackage[T1]{fontenc}
\usepackage[ngerman]{babel}
\usepackage{lmodern}
\usepackage{microtype}

\usepackage{amsmath,amssymb,amsfonts,mathtools}
\usepackage{multicol}
\usepackage{enumitem}
\usepackage{array,tabularx,booktabs,makecell}
\usepackage[dvipsnames,table]{xcolor}
\usepackage[explicit]{titlesec}
\usepackage{etoolbox}
\usepackage{fancyhdr}
\usepackage{lastpage}

% ---------- Farben ----------
\definecolor{secblue}{HTML}{1f4e79}
\definecolor{subblue}{HTML}{2e75b6}
\definecolor{boxbg}{HTML}{f2f2f2}
\definecolor{accent}{HTML}{c00000}
\definecolor{rulegray}{HTML}{888888}

% ---------- Spalten ----------
\setlength{\columnsep}{3mm}
\setlength{\columnseprule}{0.3pt}
\def\columnseprulecolor{\color{rulegray}}

% ---------- Listen kompakt ----------
\setlist{nosep,leftmargin=1.1em,labelsep=0.35em,topsep=0.2ex,partopsep=0pt}
\setlist[itemize]{label=\textbullet}

% ---------- Abstaende ----------
\setlength{\parindent}{0pt}
\setlength{\parskip}{0.15ex}
\setlength{\abovedisplayskip}{0.3ex}
\setlength{\belowdisplayskip}{0.3ex}
\setlength{\abovedisplayshortskip}{0.2ex}
\setlength{\belowdisplayshortskip}{0.2ex}
\setlength{\jot}{1pt}

% ---------- Tabellen kompakt ----------
\renewcommand{\arraystretch}{1.05}
\setlength{\tabcolsep}{3pt}
\setlength{\arrayrulewidth}{0.3pt}
\arrayrulecolor{rulegray}

% ---------- Section-Stil (kein Vollblock; dick + farbige Linie darunter) ----------
\titleformat{\section}
  {\large\bfseries\color{secblue}}
  {}{0pt}{#1}
  [\vspace{-0.4ex}{\color{secblue}\hrule height 0.8pt}\vspace{0.1ex}]
\titlespacing*{\section}{0pt}{1.4ex plus 0.4ex}{0.5ex}

\titleformat{\subsection}
  {\small\bfseries\color{subblue}}
  {}{0pt}{#1}
\titlespacing*{\subsection}{0pt}{0.9ex plus 0.2ex}{0.2ex}

\titleformat{\subsubsection}
  {\small\bfseries}
  {}{0pt}{#1}
\titlespacing*{\subsubsection}{0pt}{0.5ex}{0.1ex}

% ---------- Mini-Helfer ----------

% \fact{...} Hervorgehobene Regel (gerahmt, hellgrau)
\newcommand{\fact}[1]{%
  \begingroup\setlength{\fboxsep}{2pt}%
  \colorbox{boxbg}{\parbox{\dimexpr\linewidth-2\fboxsep}{\small #1\strut}}%
  \endgroup\par\vspace{0.2ex}}

% \warn{...} Stolperfalle (roter Strich links)
\newcommand{\warn}[1]{%
  \begingroup\setlength{\fboxsep}{2pt}%
  \noindent{\color{accent}\rule{2pt}{\baselineskip}}\hspace{2pt}%
  \begin{minipage}[t]{\dimexpr\linewidth-6pt}\small #1\end{minipage}%
  \endgroup\par\vspace{0.2ex}}

% Math-Shortcuts
\newcommand{\R}{\mathbb{R}}
\newcommand{\N}{\mathbb{N}}
\newcommand{\Z}{\mathbb{Z}}
\newcommand{\Q}{\mathbb{Q}}
\newcommand{\C}{\mathbb{C}}
\DeclareMathOperator{\rg}{rg}
\DeclareMathOperator{\spur}{tr}
\DeclareMathOperator{\im}{im}
\DeclareMathOperator{\diag}{diag}
\DeclareMathOperator{\sign}{sgn}
\DeclareMathOperator*{\argmin}{arg\,min}
\DeclareMathOperator*{\argmax}{arg\,max}
\newcommand{\trans}{^{\!\top}}
\newcommand{\norm}[1]{\left\lVert #1 \right\rVert}
\newcommand{\abs}[1]{\left\lvert #1 \right\rvert}
\newcommand{\inner}[2]{\langle #1,\,#2\rangle}
\newcommand{\dx}{\,\mathrm{d}x}
\newcommand{\dy}{\,\mathrm{d}y}
\newcommand{\dt}{\,\mathrm{d}t}

% Zeigt eine Display-Formel zentriert; skaliert nur runter wenn zu breit.
\newlength{\fitw}
\newcommand{\fitline}[1]{%
  \par\nointerlineskip
  \settowidth{\fitw}{$\displaystyle #1$}%
  \ifdim\fitw>\linewidth
    \noindent\resizebox{\linewidth}{!}{$\displaystyle #1$}\par
  \else
    \[\displaystyle #1\]%
  \fi}

% Tabular fuer fixe Spaltenbreiten (Spickzettel-typisch)
\newenvironment{tightt}[1]
  {\par\noindent\begin{tabular}{#1}}
  {\end{tabular}\par}
% Fuer X-Spalten direkt \begin{tabularx}{\linewidth}{...} verwenden
% (kann nicht in \newenvironment gewrappt werden -- TX@get@body Konflikt)

% Globale Schutzschalter gegen Overflow
\emergencystretch=2em
\tolerance=2000
\hbadness=10000
\hfuzz=2pt
\vfuzz=2pt

% ---------- Kopf-/Fusszeile ----------
\pagestyle{fancy}
\fancyhf{}
\renewcommand{\headrulewidth}{0pt}
\renewcommand{\footrulewidth}{0pt}
\fancyfoot[L]{\scriptsize\color{rulegray}\textit{\sheettitle}}
\fancyfoot[R]{\scriptsize\color{rulegray}Seite \thepage\,/\,\pageref{LastPage}}
\setlength{\footskip}{8pt}

% Titel-Header oben auf Seite 1
\newcommand{\sheetheader}[2]{%
  {\Large\bfseries\color{secblue} #1}\hfill
  {\small\color{rulegray} #2}\par
  \vspace{0.4ex}{\color{secblue}\hrule height 0.7pt}\vspace{1ex}}

% Default-Titel (wird ueberschrieben)
\newcommand{\sheettitle}{Open-Book Spickzettel}
